\section{Síntesis y ampliación}

Esta clase construyó el panorama completo de la escalación de cómputo en tiempo de prueba. Repeated sampling obtiene una escalación de coverage predecible mediante la probabilidad de éxito de cola larga; el generation--verification gap muestra que la cobertura de candidatos no se convierte automáticamente en precisión final; compute-optimal scaling mejora el rendimiento por unidad de cómputo mediante una proporción dinámica entre exploración paralela y revisión secuencial; Archon, por su parte, eleva la combinación de generator, critic, ranker, fuser y verifier a una arquitectura de sistema que puede buscarse.

Se puede entender de forma aproximada el desempeño final como:

$$
\mathrm{Delivered\ Quality}
=\mathrm{Generation\ Coverage}
\times\mathrm{Selection\ Reliability}
\times\mathrm{Budget\ Efficiency}.
$$

Si cualquiera de los tres factores se acerca a cero, el cómputo de inferencia adicional pierde su efecto. La siguiente clase se centra, por ello, en robust verification: cómo entrenar el outcome verifier y el process verifier, cómo construir automáticamente supervisión a nivel de paso, y cómo usar múltiples verifiers débiles para reducir la brecha de selección.

\end{document}
